\documentclass[11pt,letterpaper]{article}
\usepackage[margin=0.75in]{geometry}
\usepackage[hidelinks]{hyperref}

\setlength{\parindent}{0pt}
\setlength{\parskip}{0.25em}
\setlength{\emergencystretch}{1em}
\pagestyle{plain}

% Team information.
\newcommand{\MemberOne}{Bingyu Chen}
\newcommand{\MemberTwo}{Yaxin Wang}

\newcommand{\proposalheading}[1]{%
  \par\addvspace{0.55em}\noindent\textbf{#1}\par\nobreak\vspace{0.08em}}
\newenvironment{compactitems}{%
  \begin{list}{\arabic{enumi}.}{\usecounter{enumi}%
  \setlength{\itemsep}{0.18em}\setlength{\topsep}{0.2em}%
  \setlength{\partopsep}{0pt}\setlength{\parsep}{0pt}}%
}{\end{list}}

\begin{document}

\begin{center}
{\large\bfseries Project 20: Dataflow Pipelines on Vortex}\\[0.25em]
CS 254A, Fall 2026 --- Project Proposal\\[0.25em]
\MemberOne{} (withheld)\quad and\quad \MemberTwo{} (withheld)
\end{center}
\vspace{-0.45em}

\proposalheading{1. Why this project}
We chose this project because we want to learn more about GPU architecture through
hands-on experiments. Vortex lets us work with C++ simulation while also exploring
the corresponding RTL. Our interests in digital IC design, computer architecture,
and AI accelerators also make this a good opportunity to explore how AI and other
algorithm workloads are mapped to GPUs, and how their execution can be improved.

\proposalheading{2. Problem and goal}
An operator chain can expose memory stalls and materialize intermediate tensors even
when different hardware units could process different tiles concurrently. However,
splitting work into pipeline stages also consumes execution resources and introduces
synchronization and queue-management overhead. Our question is:
\emph{under what conditions can explicit producer--consumer pipelines improve
end-to-end latency and data locality on Vortex?}
We aim to adapt the core warp-specialization and asynchronous synchronization ideas
of WASP~\cite{wasp} to Vortex, then extend communication across cores using
Kitsune-inspired queues and placement~\cite{kitsune}.

\proposalheading{3. Approach and preliminary exploration}
\textbf{Starting point.}
We have built Vortex v3.0 and run the existing asynchronous-copy (DXA) and
warpgroup matrix-multiply (WGMMA) examples in SimX~\cite{vortex}.
We can collect performance counters and export Perfetto traces.
In a one-core FP16-input/FP32-output GEMM experiment with $M=N=128$ and
$K\in\{32,64,128,256\}$, double buffering ranged from 4.67\% fewer to 5.10\% more
cycles than single buffering. All cases passed the existing numerical checks.
These initial results suggest that extra buffering alone is not always helpful;
we would like to understand when overlapping stages is worth its overhead.

\textbf{Planned implementation.}
\begin{compactitems}
\item \textbf{Start with a small operator chain.}
Use data loading/staging, GEMM, and ReLU to compute $Y=\mathrm{ReLU}(AB)$.
Compare separate kernels with a conventional fused kernel, so we can see what
fusion already saves before adding a pipeline.
\item \textbf{Try different warp roles within one core.}
Adapt the C++ DXA/WGMMA example so producer warps load tiles while consumer warps
perform matrix multiplication, then add activation processing. Shared-memory
buffers and arrive/wait barriers will let the stages exchange tiles safely.
We will explore tile sizes, buffer counts, and warp allocations to see where
useful overlap appears and where stages spend time waiting.
\item \textbf{Let stages run on different cores.}
Build bounded queues using global memory and atomics, with L2 as the intended
communication path. Extend the SimX kernel-management unit (KMU) to place stages
on selected cores. We will first check that tiles arrive correctly and execution
makes progress, then study how placement and queue size affect traffic and latency.
\end{compactitems}
We will start in SimX, using WASP for warp specialization and asynchronous
coordination, and Kitsune for cross-core pipelines. Alongside the C++ work, we will
trace the relevant scheduling and synchronization paths in RTL and implement and
test corresponding RTL changes where our design changes hardware behavior.
Our report will explain which paper mechanisms we reproduced or adapted for Vortex.

\newpage

\proposalheading{4. Expected outcome and evaluation}
We hope to finish with working intra-core and cross-core load--GEMM--ReLU examples,
queue and placement support, and repeatable experiment scripts. We want to understand
when specialized stages help and when a simpler fused kernel works better.

We will compare separate kernels, conventional fusion, and both pipeline versions
using the same total hardware resources and precision. We will measure end-to-end
device cycles, throughput, and off-chip traffic, and check results against a CPU
reference. Experiments will vary matrix/tile sizes, buffer depths, and stage allocations;
counters and traces will help explain stage waiting. Success means correct execution
and a clear account of where pipelining saves time or traffic and where it slows down.

\proposalheading{5. Stretch goal}
Our main extension would be a small Transformer feed-forward block using
SwiGLU~\cite{glu}:
\[
H=\mathrm{SiLU}(XW_g)\odot(XW_u),\qquad Y=HW_d.
\]
Two GEMM branches feed an activation and elementwise multiply, followed by an
output GEMM. We would explore passing intermediate tiles between stages instead of
writing full tensors to off-chip memory, comparing separate, fused, and pipelined
kernels across input sizes. This tests whether pipelining remains useful when stages
must join two input streams.
As a further architectural exploration, DICE~\cite{dice} motivates estimating
whether direct forwarding between simple stages could save buffer accesses,
starting with a small SimX model and access counts.

\proposalheading{6. Milestones and schedule}
We plan roughly seven weeks of work, ending with submission on December~10, 2026:
\begin{compactitems}
\item \textbf{Weeks 1--2: baselines and design exploration.}
Read the papers, reproduce baselines, build separate/fused chain examples,
and trace DXA, WGMMA, and barriers through SimX and RTL.
\item \textbf{Weeks 3--4: intra-core pipeline.}
Implement loader/compute roles, tile-ready and buffer-release synchronization,
and ReLU; compare buffering and warp allocation.
\item \textbf{Weeks 5--6: cross-core pipeline.}
Implement atomic queues and KMU placement; test ordering and progress,
then study stage balance and communication cost. Attempt the stretch workload if time permits.
\item \textbf{Final week: evaluation and report.}
Finish comparisons, explain results with counters/traces, and package code,
scripts, and the report.
\end{compactitems}

\proposalheading{7. Risks and instructor feedback}
The main risk is that synchronization and communication overhead may limit the
benefits of pipelining. We would welcome the instructor's feedback on what this
project should cover and what outcomes to aim for, including any specific
requirements, suggested priorities, or ideas for further exploration.

\begingroup
\footnotesize
\setlength{\parskip}{0pt}
\begin{thebibliography}{5}
\setlength{\itemsep}{0.15em}
\bibitem{wasp}
N.~C.~Crago, S.~Damani, K.~Sankaralingam, and S.~W.~Keckler,
\emph{WASP: Exploiting GPU Pipeline Parallelism with Hardware-Accelerated Automatic
Warp Specialization}, HPCA, 2024.
\href{https://doi.org/10.1109/HPCA57654.2024.00086}{doi:10.1109/HPCA57654.2024.00086}.

\bibitem{kitsune}
M.~Davies, N.~Crago, K.~Sankaralingam, and S.~W.~Keckler,
\emph{Kitsune: Enabling Dataflow Execution on GPUs with Spatial Pipelines},
ACM TACO, vol.~22, no.~4, art.~146, 2025.
\href{https://doi.org/10.1145/3777466}{doi:10.1145/3777466}.

\bibitem{vortex}
Vortex GPGPU, \emph{Release v3.0}, source commit \texttt{7fe1936}.
\href{https://github.com/vortexgpgpu/vortex/releases/tag/v3.0}{Release notes and documentation}.

\bibitem{glu}
N.~Shazeer, \emph{GLU Variants Improve Transformer}, 2020.
\href{https://arxiv.org/abs/2002.05202}{arXiv:2002.05202}.

\bibitem{dice}
J.~Wang, A.~D.~Lu, Z.~Zeng, and A.~Li,
\emph{DICE: Enabling Efficient General-Purpose SIMT Execution with Statically
Scheduled Coarse-Grained Reconfigurable Arrays}, 2026.
\href{https://arxiv.org/abs/2605.05496}{arXiv:2605.05496}.
\end{thebibliography}
\endgroup

\end{document}
